% math-spine/collapse-risk-geometry.tex
% Shared-Horizon Preprint — Mathematical Spine
% Module 6: Collapse-Risk Geometry

\section{Collapse-Risk Geometry}
\label{sec:collapse-risk-geometry}

\subsection{Collapse-Risk Function}

Let $\mathcal{X}$ be any lineage manifold (operator, expressive, or narrative).
We define a collapse-risk function
\[
C : \mathcal{X} \rightarrow \mathbb{R}_{\ge 0}
\]
that measures the degree to which a point $p \in \mathcal{X}$ approaches
another lineage’s geometry.

The function $C$ must satisfy:

\begin{enumerate}
    \item Non-negativity: $C(p) \ge 0$.
    \item Smoothness: $C \in C^\infty(\mathcal{X})$.
    \item Minimality: $C(p) = 0$ iff $p$ is lineage-stable.
\end{enumerate}

\subsection{Collapse Boundary}

The collapse boundary is defined as:
\[
\partial \mathcal{C} = \{ p \in \mathcal{X} \mid C(p) = \delta \}
\]
for some lineage-specific threshold $\delta > 0$.

Crossing $\partial \mathcal{C}$ indicates conceptual collapse, such as:

\begin{itemize}
    \item expressive $\rightarrow$ operator collapse,
    \item narrative $\rightarrow$ expressive collapse,
    \item operator $\rightarrow$ expressive collapse,
    \item spectral $\rightarrow$ narrative collapse.
\end{itemize}

\subsection{Collapse Gradient}

The collapse gradient is defined as:
\[
\nabla C(p) \in T_p \mathcal{X}.
\]

The magnitude of the gradient indicates collapse acceleration:
\[
\| \nabla C(p) \| \quad \text{large} \quad \Rightarrow \quad \text{high collapse acceleration}.
\]

Collapse-safe transitions must satisfy:
\[
\| \nabla C(T(p)) \| \le \| \nabla C(p) \|.
\]

\subsection{Collapse Curvature}

Collapse curvature is defined as:
\[
K_C(p) = \text{Curv}_{g_{\mathcal{X}}}(C)(p).
\]

High collapse curvature indicates unstable regions where small perturbations
cause large lineage drift.

Collapse-safe regions satisfy:
\[
|K_C(p)| < K_{\max}.
\]

\subsection{Collapse-Safe Transitions}

A transition map $\Psi : \mathcal{X} \rightarrow \mathcal{X}$ is collapse-safe if:
\[
C(\Psi(p)) \le C(p)
\]
and
\[
\| \nabla C(\Psi(p)) \| \le \| \nabla C(p) \|.
\]

Collapse-safe transitions preserve lineage sovereignty.

\subsection{Collapse-Risk Tensor}

We define the collapse-risk tensor:
\[
\mathcal{T}_C(p) = \nabla^2 C(p)
\]
which measures second-order collapse behavior.

Collapse-safe regions satisfy:
\[
\| \mathcal{T}_C(p) \| < T_{\max}.
\]

This tensor is used to detect collapse-prone curvature wells.

\subsection{Cross-Lineage Collapse Geometry}

For two lineages $\mathcal{X}$ and $\mathcal{Y}$, collapse geometry is defined
using a cross-lineage collapse map:
\[
\Gamma : \mathcal{X} \rightarrow \mathcal{Y}.
\]

Collapse occurs when:
\[
C_{\mathcal{Y}}(\Gamma(p)) > \delta_{\mathcal{Y}}.
\]

Cross-lineage collapse-safe maps satisfy:
\[
C_{\mathcal{Y}}(\Gamma(p)) \le C_{\mathcal{X}}(p).
\]

\subsection{Collapse-Risk Geometry Summary}

Collapse-risk geometry defines the mathematical boundaries that prevent
conceptual collapse between lineages. It uses collapse functions, gradients,
curvature, tensors, and boundary conditions to ensure that operator, expressive,
narrative, and spectral systems remain distinct and sovereign within the
Shared-Horizon mosaic.
